%%%%%%%%%%%%%%%%%%%%%%%%%%%%%%%%%%%%%%%%%%%%%%%%%%%%%%%%%%%%%%%%%%%%%%%%%%%%
%% Author template for Management Science (mnsc) for articles with e-companion (EC)
%% Mirko Janc, Ph.D., INFORMS, mirko.janc@informs.org
%% ver. 0.95, December 2010
%%%%%%%%%%%%%%%%%%%%%%%%%%%%%%%%%%%%%%%%%%%%%%%%%%%%%%%%%%%%%%%%%%%%%%%%%%%%
%\documentclass[msom,blindrev]{informs3_no_remark} 
%\documentclass[nonblindrev]{informs3_no_remark} 
\documentclass[opre,nonblindrev]{informs4} 
% current default for manuscript submission
%\documentclass[opre,nonblindrev]{informs3}

\OneAndAHalfSpacedXI % current default line spacing
%\OneAndAHalfSpacedXII 
%\DoubleSpacedXII
%%\DoubleSpacedXI

% If hyperref is used, dvi-to-ps driver of choice must be declared as
%   an additional option to the \documentstyle. For example
%\documentclass[dvips,mnsc]{informs3}      % if dvips is used
%\documentclass[dvipsone,mnsc]{informs3}   % if dvipsone is used, etc.

% Private macros here (check that there is no clash with the style)

% Natbib setup for author-year style
\usepackage{natbib}
 \bibpunct[, ]{(}{)}{,}{a}{}{,}%
 \def\bibfont{\small}%
 \def\bibsep{\smallskipamount}%
 \def\bibhang{24pt}%
 \def\newblock{\ }%
 \def\BIBand{and}%
 
 \usepackage{tikz,calc}
 \usetikzlibrary{decorations.markings}
 \usepackage{soul}
 \setul{}{0.75pt}
 \setuloverlap{1.5pt}


\usepackage{subfigure}
%\usepackage{subcaption}
\usepackage{mwe}
\usepackage{lscape,booktabs,longtable}
\usepackage{amsmath}
\usepackage[pdftex,colorlinks=true,urlcolor=blue,citecolor=blue,pdfstartview=FitH]{hyperref}
\usepackage{comment}
\usepackage{enumerate}
\usepackage[ruled,vlined]{algorithm2e}
\usepackage[graphicx]{realboxes}

\newcommand{\modone}{\mod_{\hspace{-.04in}1\hspace{.04in}}}

\newcommand{\tr}{\textcolor{red}}
\newcommand{\tb}{\textcolor{blue}}
\newcommand{\tg}{\textcolor{mygreen}}

\newcommand{\E}[1]{{\mathrm{E}\left[#1\right]}}
\newcommand{\f}[1]{{\text{f}\left(#1\right)}}
\newcommand{\ind}[1]{{\,\mathbf{1}\hspace{-.025in}\left\{#1\right\}}}
\newcommand{\PP}[1]{\mathrm{P}\left( #1 \right)}
\newcommand{\Var}[1]{{\mathrm{Var}\left(#1\right)}}
\newcommand{\Cov}[1]{{\mathrm{Cov}[#1]}}
\newcommand{\T}{{\mathrm{T}}}

%%%% Water Mark

%\usepackage{draftwatermark}
%\SetWatermarkText{Under Review}
%\SetWatermarkScale{0.5}


%\newtheorem{hypothesis}[theorem]{Hypothesis}
%\newtheorem{hypothesis}{Hypothesis}

%% Setup of theorem styles. Outcomment only one.
%% Preferred default is the first option.
\TheoremsNumberedThrough     % Preferred (Theorem 1, Lemma 1, Theorem 2)
%\TheoremsNumberedByChapter  % (Theorem 1.1, Lema 1.1, Theorem 1.2)
\ECRepeatTheorems

%% Setup of the equation numbering system. Outcomment only one.
%% Preferred default is the first option.
\EquationsNumberedThrough    % Default: (1), (2), ...
%\EquationsNumberedBySection % (1.1), (1.2), ...

% For new submissions, leave this number blank.
% For revisions, input the manuscript number assigned by the on-line
% system along with a suffix ".Rx" where x is the revision number.
%\MANUSCRIPTNO{OPRE-2020-04-226}



\begin{document}


% For amer: \(\sum_{n=1}^\infty (\prod_{\ell=0}^{i-1}) \frac{\kappa_{\ell+1}^+ - \ell}{\kappa_{\ell+1}^+ - \ell + s} (n+1)^{n-1} \frac{\rho^n e^{-\rho(n+1)}}{n!}\)

\textcolor{red}{Amer Note: $\tau_N$ denotes the \textit{last} point in the Hawkes Process, i.e. the \textit{duration}.}

\begin{theorem}

\[
\mathbb{E}\left[e^{-s\tau_N}\right]
= \mathbb{E}\left[\mathsf{G_B( s)}\right] - e^{-\rho} \mathsf{G_1( s)}
\quad \text{where } \quad
\mathsf{G_B(s)} = \prod_{\ell=0}^{B-1} \left(1 + \frac{s}{\kappa_{\ell+1}^+ - \ell}\right)^{-1}, B \sim \mathrm{Borel}(\rho)
\]

\end{theorem}

\proof{Proof.}

By the Law of Total Expectation,  
\[
\mathbb{E}\!\left[e^{-s\tau_N}\right]
= \sum_{n=1}^{\infty} \underbrace{\mathbb{P}(N=n+1)}_{\textcircled{1}}
\, \underbrace{\mathbb{E}\!\left[e^{-s\tau_N}\mid N=n+1\right]}_{\textcircled{2}}.
\]

\textcircled{1} We know, following the proof of Theorem 1, that  
\[
\mathbb{P}(N=n+1) = (n+1)^{n-1}\frac{\rho^{n} e^{-\rho(n+1)}}{n!}
\]

\textcircled{2} We know from the theory of LSTs that if \(X \overset{d}{=} Y\), then
\(\mathbb{E}[e^{-sX}] = \mathbb{E}[e^{-sY}]\).

From Daw and Yom-Tov (2024, Thm 6), we know that  
\[
\tau_N \overset{d}{=} \sum_{\ell=0}^{N-1} Z_{\pi_\ell,\ell}
\quad \text{conditioned on } N=n+1.
\]

Then  
\[
\mathbb{E}\!\left[e^{-s\tau_N}\mid N=n+1\right]
= \mathbb{E}\!\left[e^{-s\sum_{\ell=0}^{n} Z_{\pi_\ell,\ell}}\right]
= \prod_{\ell=0}^{n} \mathbb{E}\!\left[e^{-sZ_{\pi_\ell,\ell}}\right],
\]
since the LST of a sum of random variables is the product of their individual LSTs.

Since $\exp(\lambda)$ has a well-known LST $\frac{\lambda}{\lambda-s}$ and \(Z_{\pi,\ell} \sim \exp(\kappa_{\ell+1}^+ - \ell)\) $$,  
\[
\mathbb{E}\!\left[e^{-s\tau_N}\mid N=n+1\right]
= \prod_{\ell=0}^{n}
\frac{\kappa_{\ell+1}^+ - \ell}{\kappa_{\ell+1}^+ - \ell + s}.
\]

Overall,  
\[
\mathbb{E}\!\left[e^{-s\tau_N}\right]
= \sum_{n=1}^{\infty}
(n+1)^{n-1}\frac{\rho^{n} e^{-\rho(n+1)}}{n!}
\prod_{\ell=0}^{n}
\frac{\kappa_{\ell+1}^+ - \ell}{\kappa_{\ell+1}^+ - \ell + s}.
\]

This can be written as  
\[
\mathbb{E}\!\left[e^{-s\tau_N}\right]
= \sum_{n=1}^{\infty}
\frac{(\rho(n+1))^{n} e^{-\rho(n+1)}}{(n+1)!}
\prod_{\ell=0}^{n}
\left(1+\frac{s}{\kappa_{\ell+1}^+ - \ell}\right)^{-1}.
\]

Changing indices (\(j = n+1\)) gives  
\[
\mathbb{E}\!\left[e^{-s\tau_N}\right]
= \sum_{j=2}^{\infty}
\frac{(\rho j)^{j-1} e^{-\rho j}}{j!}
\prod_{\ell=0}^{j-1}
\left(1+\frac{s}{\kappa_{\ell+1}^+ - \ell}\right)^{-1}.
\]

This can be expressed equivalently as  
\[
\mathbb{E}\!\left[e^{-s\tau_N}\right]
= \sum_{j=2}^{\infty}
\mathbb{P}(B=j)
\prod_{\ell=0}^{B-1}
\left(1+\frac{s}{\kappa_{\ell+1}^+ - \ell}\right)^{-1}.
\]

Where $B \sim \mathrm{Borel}(\rho)$

Finally, adding back the \(j=1\) case and subtracting it explicitly,
\[
\mathbb{E}\!\left[e^{-s\tau_N}\right]
= \sum_{j=1}^{\infty}
\mathbb{P}(B=j)
\prod_{\ell=0}^{B-1}
\left(1+\frac{s}{\kappa_{\ell+1}^+ - \ell}\right)^{-1}
- \mathbb{P}(B=1)
\prod_{\ell=0}^{0}
\left(1+\frac{s}{\kappa_{\ell+1}^+ - \ell}\right)^{-1}.
\]

Now, letting 
\[
\mathsf{G_x(s)}=\prod_{\ell=0}^{x-1}
\left(1+\frac{s}{\kappa_{\ell+1}^+ - \ell}\right)^{-1} 
\]

We retrieve
\[E[e^{-s\tau_N}]= \sum\limits_{j=1}^\infty \mathbb{P}(B=j) \mathsf{G_B( s)} - \mathbb{P}(B=1) \mathsf{G_1( s)} =  \mathbb{E}\!\left[\mathsf{G_B( s)}\right]
- e^{-\rho}\mathsf{G_1( s)} \]  $\square$


\textcolor{red}{Amer: Seems nice! following the style of thm 1, I tried to isolate the randomness here to \textit{just} the Borel random variable. I tried following your notation for $H_B$, showing that this is a expression depending on $B$ (but it also depends on $s$ since it's an LST which uses $s$). Also note $G_B(s)$ is monotonic in $s$, which hopefully makes it easier to approximate or think about. Also, ideally products like this play nicely with probabilities/set intersections/bounds. See what you think!}

\textcolor{red}{I had ChatGPT plot a quick example of the LST where $\kappa_{\ell+1}^+ = \ell + 2$ and $\rho=0.5$ as a sanity check. Monotonic in $s$!}

\begin{figure}[htbp]
    \centering
    \includegraphics[width=0.5\linewidth]{amers/LST_graph.png}
\end{figure}


help!

\[ \mathbb{P} ({\tau < t}) = \sum_n^\infty (\mathbb{P}( \tau < t | N=n)) \underbrace{(\mathbb{P}(N=n))}_{\text{Borel}} \]

\[ \mathbb{P} ({\tau < t}) = \sum_n^\infty \sum_{h \in H_n} (\mathbb{P} \sum_{j=1}^n \sum_{i=1}^{h_j} \text{Exp(j)} \mid h, N=n)\underbrace{(\mathbb{P}(h))}_{\text{number of d in h / total d}}\underbrace{(\mathbb{P}(N=n))}_{\text{Borel}} \]

\[ \mathbb{P} ({\tau < t}) = \sum_n^\infty \sum_{h \in H_n} (\mathbb{P} \sum_{j=1}^n \Gamma(h_j, j) \mid h, N=n)\underbrace{(\mathbb{P}(h))}_{\text{number of d in h / total d}}\underbrace{(\mathbb{P}(N=n))}_{\text{Borel}} \]

Notes:

\begin{enumerate}
    \item All Gammas are \textit{conditionally} independent given $h,n$
    \item pdf of sum of independent gammas is known (?) in explicit form (\hyperlink{https://www-personal.umd.umich.edu/~fmassey/gammaRV/}{here}). 
    \begin{enumerate}
        \item Note that the shape $h_j$ is FIXED (not random) because it is inside the $h$-conditioning. 
        \item Additionally, the rate $i$ is a deterministic index, not random.
    \end{enumerate}
    I.E. these gammas are deterministic, the pdf should be expressable in closed form.
    \item $h \in H_n$ is \textit{deterministically and exhaustively} identified by all feasible integer lattice points within the polytope:
    \begin{enumerate}
        \item $1 \le h_i^+ \le n$
        \item $ h_i^+ \ge \min(h_{i-1}^+ +1, n)$
    \end{enumerate}
    Which is affinely isomorphic (geniunely just a translation up) to the $n-1$ dimensional simplex $\Delta_n$. Moreover, $H_n = \bigsqcup_{m=1}^n H_m$ where $H_m$ can be deterministically characterized by the integer lattice points inside the polytope (translation) of $\Delta_m$ given by 
    \begin{enumerate}
        \item $1 < h_i$
        \item $ n = \sum_{i=1}^m h_i $ 
    \end{enumerate}
\end{enumerate}


\section*{Acknowledgements} 


\tb{we can't forget to thank the NSF!}

\begin{proposition}[Multinomial rotation representation]
Conditional on $N=n+1$, let
\[
(X_1,\ldots,X_{n+1})
\sim
\operatorname{Multinomial}
\left(n;\frac{1}{n+1},\ldots,\frac{1}{n+1}\right).
\]
Define
\[
S_0=0,
\qquad
S_j=\sum_{i=1}^{j}(X_i-1),
\quad 1\leq j\leq n,
\]
and let $r$ be the first index at which $S_0,\ldots,S_n$
attains its minimum. Rotate the vector so that it begins immediately
after $r$. The final coordinate of the rotated vector is zero, and,
after deleting it, the remaining vector has the distribution of
$\boldsymbol{\kappa}\mid N=n+1$.
\end{proposition}

\begin{proof}
There are $n$ balls distributed among $n+1$ coordinates, so
\[
\sum_{i=1}^{n+1}(X_i-1)=-1.
\]
The cycle lemma states that a sequence of integers, each at least
$-1$ and with total sum $-1$, has a unique cyclic rotation whose
proper partial sums are all nonnegative. This rotation is obtained
by beginning immediately after the first minimum of the original
partial sums; see \cite[Corollary~14.4]{janson2012simply}. Therefore,
if $(Y_1,\ldots,Y_{n+1})$ denotes the rotated vector, then
\[
\sum_{i=1}^{j}Y_i\geq j,
\qquad 1\leq j\leq n.
\]
Because the total sum is $n$, the last coordinate must be $Y_{n+1}=0$.
Thus, after deleting this coordinate, the vector
$\mathbf{k}=(Y_1,\ldots,Y_n)$ satisfies
\[
k_i^+=\sum_{j=1}^{i}k_j\geq i,
\qquad 1\leq i\leq n,
\]
which is exactly the feasibility condition for the response-count
vector.

Every feasible vector $\mathbf{k}$ arises from each of the $n+1$
cyclic rotations of $(\mathbf{k},0)$. Each such multinomial vector
has probability
\[
\frac{n!}{(n+1)^n\prod_{i=1}^{n}k_i!}.
\]
It follows that
\[
\Pr(\boldsymbol{\kappa}=\mathbf{k}\mid N=n+1)
=
(n+1)\frac{n!}{(n+1)^n\prod_{i=1}^{n}k_i!}
=
\frac{n!}{(n+1)^{n-1}\prod_{i=1}^{n}k_i!},
\]
which is the conditional distribution of $\boldsymbol{\kappa}$.
\end{proof}

% BIB:
\bibliographystyle{informs2014} 
\bibliography{Bibliography.bib} 





%\begin{APPENDIX}{}
%
%\end{APPENDIX}




\end{document}